%\documentclass[12pt,letterpaper]{article}



%%%
%% Required packages:
%%
\usepackage{etex}
\usepackage{amsmath}
\usepackage{amssymb}
\usepackage{amstext}
\usepackage{amsthm}
\usepackage{fancyhdr,fancybox}
\usepackage{mathrsfs} % need for \mathscr command
\usepackage{multicol}
\usepackage[normalem]{ulem}
%\usepackage{pictex,pictexplus}
% \usepackage{pictexwd}
\usepackage{rotating}
\usepackage{url}
\usepackage{xfrac}
\usepackage{booktabs}
\usepackage{color,soul}
\usepackage{comment}

\usepackage{hyperref}
\hypersetup{
    colorlinks=true,     % false: boxed links; true: colored links
    linkcolor=blue,      % color of internal links
    citecolor=blue,      % color of links to bibliography
    filecolor=blue,      % color of file links
    urlcolor=blue        % color of external links
}


\usepackage{tikz}
\usetikzlibrary{backgrounds,calc}

%%
%% Common formatting commands
%%
\setlength{\parindent}{0in}
\setlength{\textwidth}{7in}
\setlength{\evensidemargin}{-0.25in}
\setlength{\oddsidemargin}{-0.25in}
\setlength{\parskip}{.5\baselineskip}
\setlength{\topmargin}{-0.5in}
\setlength{\textheight}{9in}

%               Problem and Part
\newcounter{problemnumber}
\newcounter{partnumber}
\newcounter{subpartnumber}
\newcommand{\Problem}{%
  \refstepcounter{problemnumber}%
  \setcounter{partnumber}{0}%
  \setcounter{subpartnumber}{0}%
  \item[\makebox{\hfil\textbf{\theproblemnumber.}\hfil}]%
  }
\newcommand{\InProblem}[2][1.6]{%
  \refstepcounter{problemnumber}%
  \setcounter{partnumber}{0}%
  \setcounter{subpartnumber}{0}%
  \textbf{\theproblemnumber.}\ \ \parbox[t]{#1in}{#2}%
}
\newcommand{\InSmallProblem}[1]{\InProblem[1.05]{#1}}
\newcommand{\NoProblem}[1]{%
  \mbox{\hfil\textbf{#1}\hfil}%
  }
\newcommand{\UnProblem}{%
  \refstepcounter{problemnumber}%
  \setcounter{partnumber}{0}%
  \setcounter{subpartnumber}{0}%
  \mbox{\hfil\textbf{\theproblemnumber.}\hfil}%
  }
\newcommand{\InFlexProblem}[1]{%
  \refstepcounter{problemnumber}%
  \setcounter{partnumber}{0}%
  \setcounter{subpartnumber}{0}%
  \textbf{\theproblemnumber.}\ \ #1%
}
%%  Part:
\newcommand{\Part}{%
  \renewcommand{\thepartnumber}{(\alph{partnumber})}%
  \refstepcounter{partnumber}
  \setcounter{subpartnumber}{0}%
  \item[(\alph{partnumber})]%
}
\newcommand{\InPart}[2][1.75]{%
  \renewcommand{\thepartnumber}{(\alph{partnumber})}%
  \refstepcounter{partnumber}%
  \setcounter{subpartnumber}{0}%
  (\alph{partnumber})\ \ \parbox[t]{#1in}{#2}%
}
\newcommand{\UnPart}{%
  \refstepcounter{partnumber}%
  \renewcommand{\thepartnumber}{(\alph{partnumber})}%
  \setcounter{subpartnumber}{0}%
  (\alph{partnumber})%
}
\newcommand{\InLargePart}[1]{\InPart[3.05]{#1}}
\newcommand{\InFullPart}[1]{\InPart[6.2]{#1}}
\newcommand{\InSmallPart}[1]{\InPart[1.05]{#1}}
%% SubPart:
\newcommand{\SubPart}{%
  \refstepcounter{subpartnumber}%
  \item[(\roman{subpartnumber})]%
  }
\newcommand{\InSubPart}[2][2.25]{%
  \stepcounter{subpartnumber}%
  (\roman{subpartnumber})\ \ \parbox[t]{#1in}{#2}%
}
\newcommand{\InSmallSubPart}[1]{\InSubPart[1.5]{#1}}
\newcommand{\InTinySubPart}[1]{%
  \stepcounter{subpartnumber}%
  (\roman{subpartnumber})\ \ \makebox{#1}%
}
\newcommand{\InMySubPart}[2][2.25]{%
  \stepcounter{subpartnumber}%
  \parbox[t]{#1in}{\makebox[0.3in][l]{(\roman{subpartnumber})}\ \ {#2}}%
}
\newcommand{\TrueFalse}{\textbf{T}\ \ \textbf{F}\ \ }
\newcommand{\TFPart}[1]{% 
  \stepcounter{partnumber}%
  \setcounter{subpartnumber}{0}%
  \item[(\alph{partnumber})] \TrueFalse\parbox[t]{5.5in}{#1}}

\newcommand{\Points}[1]{(#1 points) \ }
\newcommand{\Point}[1]{(#1 point) \ }


%% For Answers:
\newcounter{answernumber}
\newcommand{\Answer}{%
  \refstepcounter{answernumber}%
  \setcounter{partnumber}{0}%
  \setcounter{subpartnumber}{0}%
  \item[\makebox{\hfil\textbf{\theanswernumber.}\hfil}]%
  }


% Setting up the headers for the first page
%\chead{} % will default to what is typed in the file.
\fancyhead[L]{\textsc{Math 22a}}
\fancyhead[R]{\textsc{Fall 2024}}
\fancyfoot[R]{
	\vspace*{\fill}\footnotesize\textcopyright{} \emph{Harvard Math 22a}	
}
\renewcommand{\headrulewidth}{0pt}

%\fancyhead[C]{}
%	\chead{}	
%	\lhead{}
%	\rhead{}
%	\cfoot{}


% Putting copyright on the footer of each page
%\fancypagestyle{secondpage}{
%	\fancyhead[C]{TESTing again} % change this to be blank
%		%\chead{}	
%	\fancyhead[L]{bears} % change this to be blank
%	\fancyhead[R]{dogs} % change this to be blank
%	\fancyfoot[R]{ %keeps the footer the same
%		\vspace*{\fill}\footnotesize\textcopyright{} \emph{Harvard Math 22a}	
%	}
%}
%\renewcommand{\headrulewidth}{0pt}
%\pagestyle{fancy}
\pagestyle{empty}


%\lhead{}
%\rhead{}
%\rfoot{\vspace*{\fill}\footnotesize\textcopyright{} \emph{Harvard Math 22a}}
%\renewcommand{\headrulewidth}{0pt}
%\pagestyle{fancy}




%%
%% Common shortcut commands:
%%
\newcommand{\ds}{\displaystyle}
\newcommand{\sgn}{\operatorname{sgn}}
\renewcommand{\Pr}{\operatorname{P}}
\newcommand{\CPr}[3][{}]{\Pr_{\text{#1}}\left(#2\,|\,#3\right)}

\newcommand{\C}{\mathbb{C}}
\newcommand{\R}{\mathbb{R}}
\newcommand{\Z}{\mathbb{Z}}
\newcommand{\N}{\mathbb{N}}
\newcommand{\Q}{\mathbb{Q}}
\newcommand{\D}{\mathbb{D}}

\newcommand{\rref}{\operatorname{rref}}
\newcommand{\im}{\operatorname{im}}
\newcommand{\rank}{\operatorname{rank}}
\newcommand{\diag}{\operatorname{diag}}
\newcommand{\Span}{\operatorname{span}}
%\renewcommand{\vec}[1]{\mathbf{#1}}
\newcommand{\charpoly}{\mathscr{P}}
\newcommand{\nicefrac}[2]{\sfrac{#1}{#2}} % using xfrac package
\newcommand{\topology}{\mathcal{T}}
\newcommand{\Inn}{\operatorname{Inn}}

\newcommand{\blankbox}[2]{\fbox{\rule{#1}{0in}\rule{0in}{#2}}}

\newenvironment{myindentpar}[1]%
 {\begin{list}{}%
         {\setlength{\leftmargin}{#1}
          \setlength{\rightmargin}{\leftmargin}}%
         \item[]%
 }
 {\end{list}}
\newtheorem{theorem}{Theorem}
\newtheorem*{NStheorem}{Theorem}
\newtheorem{cor}[theorem]{Corollary}
\newtheorem*{claim}{Claim}
\newtheorem{defn}{Definition}

\newenvironment{amatrix}[1]{%
  \left(\begin{array}{@{}*{#1}{c}|c@{}}
}{%
  \end{array}\right)
}

%--------grstep
% For denoting a Gauss' reduction step.
% Use as: \grstep{\rho_1+\rho_3} or \grstep[2\rho_5 \\ 3\rho_6]{\rho_1+\rho_3}
\newcommand{\grstep}[2][\relax]{%
   \ensuremath{\mathrel{
       {\mathop{\longrightarrow}\limits^{#2\mathstrut}_{
                                     \begin{subarray}{l} #1 \end{subarray}}}}}}
\newcommand{\swap}{\leftrightarrow}


%\usetikzlibrary{arrows}
%\usepackage{wrapfig}
%
%\makeatletter
%\renewcommand*\env@matrix[1][*\c@MaxMatrixCols c]{%
	%\hskip -\arraycolsep
	%\let\@ifnextchar\new@ifnextchar
	%\array{#1}}
%\makeatother
%
%\def\env@matrix{\hskip -\arraycolsep
	%\let\@ifnextchar\new@ifnextchar
	%\array{*\c@MaxMatrixCols c}}

\section{{Dimension and Rank, $\vec\S$4.5}}% 
\chead{\Large\textbf{Dimension and Rank, $\vec\S$4.5}}% 

%\begin{document}
\thispagestyle{fancy}

Recall from last class:

\begin{itemize}
\item {\bf Theorem:} Let $\mathcal{B} =\{b_1, \dots, b_n\}$ be a basis for a  vector space $V$. Then for each $x \in V$ there exists unique scalars $c_1,\dots, c_n$ such that $x=c_1 b_1+\dots+ c_n b_n$.
\item Suppose $\mathcal{B}=\{b_1, \dots, b_n\}$ is a basis for a vector space $V$. The {\bf coordinates of $x$ relative to $\mathcal{B}$} are the weights $c_1, \dots, c_n$ such that $x=c_1 b_1 +\dots c_n b_n$. We define the {\bf coordinate mapping} from $V$ to $\mathbb{R}^n$ by $[x]_\mathcal{B}=\begin{bmatrix} c_1\\ \vdots\\ c_n \end{bmatrix}$.
\item Let $\mathcal{B}=\{b_1, \dots, b_n\}$ be a basis for $\mathbb{R}^n$, then define $P_\mathcal{B}=[b_1 \; \dots \; b_n]$. We call $P_\mathcal{B}$ the {\bf change of coordinates matrix}. Note $x = P_\mathcal{B} [x]_{\mathcal{B}}$ and $[x]_\mathcal{B}=P_{\mathcal{B}}^{-1} x.$
\item {\bf Theorem:} Let $\mathcal{B}=\{ b_1, \dots, b_n\}$ be a basis for a vector space $V$. Then the coordinate mapping $x \mapsto [x]_\mathcal{B}$ is a one-to-one, onto, linear transformation from $V$ to $\mathbb{R}^n$.
\end{itemize}

\newpage

\begin{defn}
	Let $T: V \to W$ be a linear map of vector spaces. If $T$ is one-to-one and onto, then $T$ is a {\bf vector space isomorphism}, and we say $V$ and $W$ are isomorphic. We write $V \cong W$.
\end{defn}

\begin{enumerate}

\Problem Show that $\mathbb{P}_3$ and $\mathbb{R}^4$ are isomorphic.

\vfill

{\bf Theorem:} If a vector space $V$ has a basis $\mathcal{B}=\{b_1, \dots, b_n\}$, then any set in $V$ containing more than $n$ vectors is linearly dependent.

\Problem Prove the theorem.

\vfill

{\bf Cor:} If $V$ is a vector space with basis $\mathcal{B}=\{b_1, \dots, b_n\}$, then each set of linearly independent vectors has at most $n$ elements.

{\bf Theorem:} If a vector space $V$ has a basis with $n$ elements, then every basis of $V$ must have exactly $n$ vectors.

\Problem Prove the theorem.

\vfill

\newpage

\begin{defn}
If $V$ is spanned by a finite set, then $V$ is called {\bf finite dimensional,} and the {\bf dimension} of $V$, written $\text{dim}(V)$, is the number of vectors in a basis of $V$. Define $\text{dim}(\{0_V\})=0$. If $V$ is not spanned by any finite set, then $V$ is called {\bf infinite dimensional}.
\end{defn}




\Problem Find the dimensions of the following vector spaces.

\Part $\mathbb{R}^n$

\vfill

\Part $\mathbb{P}_n$

\vfill

\Part Let $\mathbb{P}$ be the vector space of polynomials.

\vfill

\Part $H=\left\{\begin{bmatrix} s-2t\\ s+t\\ 3t\\ \end{bmatrix} : s,t \in \mathbb{R} \right\}$.

\vfill

\Part $H=\text{Span}\left\{ \begin{bmatrix} 1\\ 1\\ 0\\ \end{bmatrix}, \begin{bmatrix} -2\\ 1\\ 1\\ \end{bmatrix}\right\}$.

\vfill

\Part $H=\{ A \in M_{n \times n} : A^{T}=A \}$.

\vfill

\newpage

{\bf Theorem:} Let $H$ be a subspace of a finite dimensional vector space $V$. Any linearly independent set in $H$ can be expanded to a basis for $H$, and $H$ is finite dimensional and $\text{dim}(H)\leq \text{dim}(V)$.

\Problem Prove the theorem.

\vfill

{\bf Basis Theorem:} Let $V$ be an $n$-dimensional vector space with $n\geq 1$. Any linearly independent set of exactly $n$ vectors in $V$ is a basis for $V$. Any set of exactly $n$ vectors that spans $V$ is a basis for $V$.

\Problem Prove the theorem.

\vfill

\newpage

\begin{defn} Let $A$ be an $m\times n$ matrix with rows $a_1, \dots, a_m$. Define the {\bf row space} of $A$ by $$\text{Row}(A)=\text{Span}\{a_1, \dots, a_m\}.$$
\end{defn}

{\bf Theorem:} If $A$ is row equivalent to $B$, then $\text{Row}(A)=\text{Row}(B)$. Furthermore, if $B$ is in reduced row echelon form, the non-zero rows of $B$ form a basis for $\text{Row}(A)$ and $\text{Row}(B)$.

\Problem Prove the theorem.

\vfill

\Problem Find the dimensions of $N(A)$, $\text{Col}(A)$, and $\text{Row}(A)$ where $$A=\begin{bmatrix}
1 & -6 & 9 &0 &-2 & 0\\
0 & 1 & 2 & -4 & 5 &0\\
0 & 0 &0 & 5 & 1 &0\\
0 & 0 &0 &0 &0 &0\\
\end{bmatrix}$$

\vfill

\begin{defn}
The {\bf rank} of the $m\times n$ matrix $A$ is the dimension of $\text{Col}(A)$. We write $\text{rank}(A)=\text{dim}(\text{Col}(A))$. The {\bf nullity} of $A$ is the dimension of $N(A)$. 
\end{defn}

\newpage

{\bf Rank-Nullity Theorem:} Let $A$ be an $m\times n$ matrix.
\begin{enumerate}
\item $$\text{dim} \;\text{Row}(A)= \text{dim} \;\text{Col}(A)$$
\item $$\text{rank(A)}+\text{dim} \;N(A) =n$$
\end{enumerate}

\Problem Prove the theorem

\vfill

\vfill

\vfill

\Problem 

\Part Let $A$ be a 7 by 5 matrix. What is the largest possible rank?

\vfill

\Part Let $A$ be a 3 by 7 matrix. What is the smallest possible nullity?

\vfill

\Part Let $A$ be a 5 by 6 matrix. Suppose the solutions to $Ax=0$ are all multiple of one non-zero vector. Will the matrix equation $Ax=b$ always be consistent?

\vfill

\newpage

{\bf Invertible Matrix Theorem Continued} The following are equivalent to $A$ being invertible:
\begin{itemize}
\item[m.] The columns of $A$ form a basis of $\mathbb{R}^n$.
\item[n.] $\text{Col}(A)=\mathbb{R}^n$.
\item[o.] $\text{dim} \; \text{Col}(A)=n$.
\item[p.] $\text{rank}(A)=n$.
\item[q.] $N(A)=\{0\}$.
\item[r.] $\text{dim}\; N(A) =0$.
\end{itemize}

\Problem Prove the theorem.

\vfill

{\bf Theorem:} Let $T: V \to W$ be a linear transformation between vector spaces. If $V$ is finite dimensional, then $$\text{dim} \; \text{im}(T) + \text{dim} \; \text{ker}(T) =\text{dim}(V).$$

\end{enumerate}

\begin{comment}
\newpage
\chead{\Large\textbf{Dimension and Rank -- Answers / Solutions}}%
\lhead{}
\rhead{}
\thispagestyle{fancy}

\begin{enumerate}
  \Answer 
\begin{enumerate}
\item Since $e_1, \dots, e_n$ is a basis, $\text{dim}(\mathbb{R}^n)=n$.
\item Since $1, t, \dots, t^n$ is a basis, $\text{dim}(\mathbb{P}_n)=n+1$.
\item The set of monomials $\{t^n: n \in \mathbb{N}\cup \{0\}\}$ forms an infinite set of linearly independent vectors, and hence $\mathbb{P}$ is infinite dimensional. 
\item Since $\left\{ \begin{bmatrix} 1\\ 1\\ 0\\ \end{bmatrix}, \begin{bmatrix} -2\\ 1\\ 3\\ \end{bmatrix} \right\}$ forms a basis for $H$, we get $\text{dim}(H)=2$.
\item Since $\left\{ \begin{bmatrix} 1\\ 1\\ 0\\ \end{bmatrix}, \begin{bmatrix} -2\\ 1\\ 1\\ \end{bmatrix}\right\}$ forms a basis for $H$, we get $\text{dim}(H)=2$.
\item $H$ is the subspace of symmetric matrices. Since elements $H$ must be symmetric matrices, the elements across the diagonal must be equal. Thus once we know the elements above the diagonal, we also know the elements below the diagonal. Thus we just need to count the number of entries below the diagonal and including the diagonal. Thus the $\text{dim}(H)=1+2+3+4+\dots+n=\frac{n(n+1)}{2}$.
\end{enumerate}
\Answer
\begin{proof} If $H=\{0\}$, then $\text{dim}(H)=0\leq \text{dim}(V)$. Assume $H\neq \{0\}$ and let $S=\{s_1, \dots, s_k$ be a linearly independent set in $H$. If $S$ spans $H$, then $S$ is a basis for $H$. If not, then there exists $s_{k+1} \in H\setminus \text{Span} \; S$. We previously proved that if $s_{k+1}$ can't be written as a linearly combination of $s_1, \dots, s_k$, then $s_1, \dots, s_{k+1}$ are linearly independent. Repeat until $s_1, \dots, s_{n+k}$ spans $H$. Now we proved that the number of linearly independent vectors in $V$ can be at most $\text{dim}(V)$, so $n+k\leq \text{dim}(V)$. So eventually we get a basis and $\text{dim}(H)\leq \text{dim}(V)$.
\end{proof}

\Answer
\begin{proof} By the previous theorem, any collection of $n$ vectors in $V$ can be expanded to a basis for $V$. Since $V$ is $n$-dimensional, we can't add any linearly independent vectors to our collection, so they must already be a basis. Suppose $S$ is a set of $n$ vectors that spans $V$. Since $V$ is at least one-dimensional, we know the span of $S$ is non-zero. The Spanning Set Theorem tells us that some subset of $S$ is a basis for $V$. Since $V$ is $n$-dimensional, any basis must contain at least $n$ vectors, and hence $S$ must be a basis. 
\end{proof}

\Answer
\begin{proof} Suppose $A$ is row-equivalent to $B$, then the rows of $B$ are obtained from $A$ by linear combinations of the rows of $A$. Thus any linear combination of the rows of $B$ is a linear combination of the rows of $A$ and vice versa. Thus $\text{Row}(A)=\text{Row}(B)$.

If $B$ is in echelon form, then the non-zero rows are linearly independent (since no non-zero row is a linear combination of non-zero rows). Since the row spaces of $A$ and $B$ are the same, we get the non-zero rows of $B$ form a basis for the row space of $A$.
\end{proof}

\Answer Since there are 3 pivots, we get the dimension of the column space is 3. By the previous theorem, the dimension of the row space is also 3. Since there are 3 free variables, the dimension of the null space is 3.

\Answer
\begin{proof} 
{\bf (a)} We proved that a basis for the column space is the pivot columns of $A$, and hence the dimension of the column space is the number of pivots. Since the pivot rows of $A$ form a basis for the row space of $A$, we get $$\text{dim}(\text{Row}(A))=\text{dim}(\text{Col}(A))=\text{rank}(A).$$

{\bf (b)} We have observed that the dimension of the null space of $A$ is the number of free variables in $Ax=0$. Thus the dimension of the null space of $A$ is the number of non-pivot columns of $A$. Hence
\begin{align*}
\text{total number of columns} &= \text{non-pivot columns } + \text{pivot columns}\\
n &= \text{dim}(N(A))+\text{rank}(A).
\end{align*}
\end{proof}

\Answer
\begin{enumerate}
\item Let $A$ be a 7 by 5 matrix. Then $\text{rank}(A)=5- \text{dim}(N(A))$. Thus $\text{rank}(A)\leq 5$.
\item Let $A$ be a 3 by 7 matrix. Then $\text{dim}(N(A))=7-\text{rank}(A)\geq 7-3=4$.
\item Let $A$ bea 5 by 6 matrix and $N(A)= \text{Span}\{v\}$ for some non-zero vector $v$. Then we know the nullity is one, so by the rank nullity theorem, we get the rank must be 5. Since the column space is a subspace of $\mathbb{R}^5$, we must have $\text{Col}(A)=\mathbb{R}^5$. Thus for each $b \in \mathbb{R}^5$, the matrix equation $Ax=b$ is consistent. 
\end{enumerate}

\Answer
\begin{proof} 
$(m) \implies (n)$ by definition of basis.

$(n) \implies (o)$ by definition of dimension and Basis theorem. 

$(o) \implies (p)$ by definition of rank.

$(p) \implies (q)$ by rank-nullity theorem.

$(q) \implies (r)$ by definition of dimension.

$(r) \implies (q)$ by definition of dimension. 

$(q) \implies (d)$ Since $N(A)=\{0\}$ means $Ax=0$ has only the trivial solution, which was property (d).

$(d) \implies (m)$ We have already proved $(d)$ implies the columns of $A$ span $\mathbb{R}^n$ and the columns are linearly independent. Thus the column space of $A$ is $\mathbb{R}^n$. 
\end{proof}
\end{enumerate}  

\end{comment}
%\end{document}


%%% Local ables: 
%%% mode: latex
%%% TeX-master: t
%%% End: 
